\documentclass[12pt]{article}
\usepackage[utf8]{inputenc}
\usepackage[T1]{fontenc}
\usepackage{hyperref}
\title{Syuba Language Resources for K-12 Curriculum Development}
\author{Compiled by Claude Sonnet 5.5 (claude-sonnet-5-5) for Andr\'as Kornai}
\date{}
\begin{document}
\maketitle

\section{Introduction}
Syuba (also called Kagate, an exonym meaning ``papermaker'' that many speakers find pejorative) is a Tibetic language of Nepal, spoken mainly in Ramechhap district. It is most closely related to Yolmo. Search results gave a speaker count of about 1,500, well below the 10,000 in the task description, so the count is uncertain. The same results say the language is still spoken across generations. Very little published material exists, and nothing found is a school curriculum.

\section{Primary Educational Resources}
No primers, textbooks or school programmes for Syuba turned up in the searches. The language is not known to be taught in schools.

\section{Linguistic and Reference Materials}
\begin{itemize}
\item Lauren Gawne, ``Language Contexts: Syuba, also known as Kagate (Nepal)'', \emph{Language Documentation and Description} 13 (2017), pp.\ 65--93. It covers history, the relationship to Yolmo, language use and the literature on Syuba. \url{https://lddjournal.org/article/id/1170}
\item Wikipedia, Kagate language: a general overview that includes phonology. \url{https://en.wikipedia.org/wiki/Kagate_language}
\item Wikipedia, Yolmo language: a grammatical overview of the closely related Yolmo. \url{https://en.wikipedia.org/wiki/Yolmo_language}
\item Wikibooks, Research on Hyolmo: related material on Hyolmo. \url{https://en.wikibooks.org/wiki/Research_on_Hyolmo}
\end{itemize}
A search summary mentions a Syuba--Nepali--English dictionary with a modified Devanagari orthography and a Roman orthography that marks low tone. I saw no URL for it, so it is not listed with a link. Gawne's sketch grammar and dictionary cover Lamjung Yolmo, not Syuba.

\section{Religious and Community Resources}
No religious or community-organisation resources for Syuba were found in the searches.

\section{Digital and Audio Resources}
\begin{itemize}
\item Endangered Languages Project entry. \url{https://endangeredlanguages.com/node/112115}
\end{itemize}
No audio or video collections were found.

\section{Curriculum Development Context}
Syuba is spoken by all age groups, which helps mother-tongue instruction. Published teaching material is almost absent. A curriculum would need new primers and a standard orthography. It could build on the Syuba dictionary orthography and on Gawne's documentation, with Yolmo materials as a comparative aid. The community's preference for the name Syuba should be respected.
\end{document}